\chapter{基因不是一直开着的灯}

\begin{kaichang}
想象你不小心被纸划破了手指。过几天，伤口长好了——长回来的是皮肤，不多也不少。不会长出一片指甲，不会冒出一截软骨，更不会长出一块肝。可是，伤口旁边的每一个细胞，细胞核里都装着同一套完整的 DNA：制造肝细胞的基因在，制造眼角膜的基因在，制造指甲的基因也在。

再看几个同样古怪的事实。萤火虫的屁股会发光，而你臀部的细胞里并没有“发光基因”坏掉了这一说——它们压根不需要那套图纸；一株月季剪下一根枝条插进土里，枝条上的细胞本来只会做叶子和茎，却忽然改行长出根来；你锻炼三个月，胳膊变粗了，肌肉细胞并没有增加新的基因，只是同一批基因干活的分量变了。

同一本书，不同的细胞读出了完全不同的内容。它们是怎么决定读哪几页的？更重要的是：\textbf{谁}来按这个开关？先猜一个答案。这一章结束时你会发现，答案既比“有个总司令在指挥”简单，又比它精妙得多。
\end{kaichang}

\section{同一本书，不同的读法}

先把能看见的现象摆出来。

你的身体由大约两百多种细胞组成：神经元伸出长长的电线，肌肉细胞里塞满收缩的纤维，红细胞干脆把细胞核都扔了，腾出地方装血红蛋白。它们的模样、工作、寿命天差地别，但除了极少数例外，每个细胞核里的 DNA 序列都一模一样——这套序列就是第 66 章讲过的那本三十亿字母的书。

再看单细胞的世界，现象更直接。大肠杆菌这种住在你肠道里的细菌，能吃葡萄糖，也能吃乳糖（牛奶里的那种糖）。把大肠杆菌放进两种糖都有的培养液里，它会先专心致志地吃葡萄糖，对身边的乳糖理都不理；等葡萄糖吃光了，它会“愣”上几十分钟，然后才开始吃乳糖。这停顿的几十分钟里，细菌内部正在发生一场悄悄的改组——这是本章的主角登场前的幕间休息。

类似的“换档”到处都有：

\begin{itemize}[itemsep=2pt]
\item 酵母掉进面团里，氧气足的时候用呼吸作用“细嚼慢咽”葡萄糖，氧气耗尽后改走发酵路线“囫囵吞枣”，顺便产出让面团膨胀的二氧化碳（第 52、53 章讲过这两条代谢路线）。
\item 你的皮肤晒太阳后变黑，是紫外线把黑色素细胞里一批“制造黑色素”的基因的产量调高了。
\item 伤口结痂时，边缘的细胞接到信号，把“分裂增殖”相关的基因开足马力，长好之后又及时刹车。
\item 秋天的叶子变黄，是叶片细胞关掉了叶绿素的维护生产线，把有用的氮、镁拆下来回收进枝干过冬。
\end{itemize}

这些现象背后是同一件事：\textbf{基因的表达量是可以调节的}。“表达”这个词在第 68、69 章已经见过——指一段 DNA 被转录成 RNA、再翻译成蛋白质、真正干起活来的全过程。拥有某个基因，和正在使用某个基因，是两回事。

\section{开关装在分子上：乳糖操纵子}

现在把镜头推进到大肠杆菌的细胞里，看那个“愣住的几十分钟”到底在发生什么。

回忆第 68 章：转录从\textbf{启动子}开始——它是 DNA 上的一段特殊序列，像停车位一样供 RNA 聚合酶停靠。聚合酶停稳了，就顺着 DNA 往前开，把一路上的基因抄成 RNA。

大肠杆菌用来消化乳糖的机器主要是两种蛋白质：一个负责把乳糖运进细胞（转运蛋白），一个负责把乳糖切开（乳糖酶）。巧的是，制造这两种蛋白的基因在 DNA 上肩并肩排在一起，共用同一个启动子——就像几台机器共用同一个总闸。在启动子和基因之间，还夹着一小段特殊的 DNA，叫\textbf{操纵序列}。这一整套“一个总闸带一串基因”的装置，叫做\textbf{乳糖操纵子}（操纵子就是“操作单元”的意思）。

\begin{center}
\begin{tikzpicture}[x=1cm,y=1cm]
  \draw[thick] (0,0) -- (9.4,0);
  \draw[fill=gray!25] (0.4,-0.32) rectangle (1.6,0.32);
  \node at (1.0,0) {\small 启动子};
  \draw[fill=gray!25] (1.9,-0.32) rectangle (2.9,0.32);
  \node at (2.4,0) {\small 操纵序列};
  \draw[fill=blue!12] (3.2,-0.32) rectangle (5.0,0.32);
  \node at (4.1,0) {\small 乳糖酶基因};
  \draw[fill=blue!12] (5.3,-0.32) rectangle (7.0,0.32);
  \node at (6.15,0) {\small 转运蛋白基因};
  \draw[fill=blue!12] (7.3,-0.32) rectangle (9.1,0.32);
  \node at (8.2,0) {\small 第三个基因};
  \draw[fill=red!20,rounded corners=2pt] (1.95,0.42) rectangle (2.85,1.1);
  \node at (2.4,0.76) {\small 阻遏蛋白};
  \node at (4.7,-0.95) {\small 阻遏蛋白坐在操纵序列上时，RNA 聚合酶被堵住去路};
\end{tikzpicture}
\end{center}

上图是人为的表示法：方框的大小和颜色不代表真实形状，真实分子是扭成一团的蛋白质和被缠绕的 DNA。

开关的关键是一个叫\textbf{阻遏蛋白}的分子。它的形状恰好能紧紧抱住操纵序列那段 DNA。当它坐在那里时，就像停车位上横了一辆路障车，RNA 聚合酶过不去，后面一整串基因都抄不出来——灯是关着的。

那乳糖来了怎么办？乳糖进入细胞后，有一小部分会被改造成一种形状略有不同的分子（异乳糖），它恰好能卡进阻遏蛋白的一个口袋里。这一卡，阻遏蛋白的形状就变了，抱不住 DNA 了，从操纵序列上掉下来。路障撤走，RNA 聚合酶长驱直入，转运蛋白和乳糖酶被成批制造出来——灯开了。等乳糖吃完，没有分子去卡阻遏蛋白的口袋，它重新坐回 DNA 上，灯又关了。

注意这条因果链的味道：没有谁在“下命令”。乳糖分子和阻遏蛋白的结合靠的是第 22 章讲过的分子间作用力，形状互补就吸住，浓度够高就大概率吸住；结合引起蛋白质形状改变，又是第 48 章讲过的“蛋白质的形状就是它的功能”。“决定”不是谁做的，是从物理规律里长出来的。

至于“先吃葡萄糖”的挑食：葡萄糖充足时，细胞里还有第二道闸——一种能感应能量状态的信号分子浓度很低，导致这个启动子本来就很难被聚合酶找到。两道闸都放行，灯才最亮。一道信号说“有乳糖吗”，另一道说“缺能量吗”，两个条件做了一次“与”运算——你的手机芯片里满是这种逻辑门，而细菌早用了几十亿年。

\section{不是开关，是调光旋钮}

“开”和“关”是方便的说法，但请你把它当成简称。真实的基因调节更像一个调光旋钮，原因有两层。

第一层：\textbf{分子结合是概率事件}。阻遏蛋白不是焊死在 DNA 上的，它不停地结合、掉下来、再结合。乳糖浓度越高，它口袋里卡着诱导物的时间比例就越大，坐在 DNA 上的时间比例就越小。于是转录不是“有或无”，而是“多或少”——乳糖多一点，每小时抄出的 RNA 就多一点，造出的酶就多一点。

第二层：\textbf{信号的强度本身就是信息}。真核细胞里，一个基因旁边往往不止一个“停车位”，而是散布着好几段调控序列：紧挨启动子的叫调控区，远在上万字母之外、却能折回来帮忙的叫\textbf{增强子}，起反作用的叫\textbf{沉默子}。能识别这些序列的蛋白质统称\textbf{转录因子}——它们有的帮聚合酶停车，有的挡路，有的把 DNA 折一个弯让远处的增强子凑到启动子跟前。一个基因听好几个转录因子的话，一个转录因子又管好几百个基因：这就是\textbf{组合调控}。几百种转录因子排列组合，理论上能调出天文数字种“亮度模式”，这就是为什么两万来个基因够指挥两百多种细胞。

\begin{qiaoliang}
这个“旋钮”的调节范围到底有多大？拿乳糖酶来算。研究人员测过：培养液里完全没有乳糖时，一个大肠杆菌里平均只有不到 5 个乳糖酶分子——注意，是“个”，不是“摩尔”。加入足量诱导物后，同一个细胞里的乳糖酶可以涨到几千个，占细胞全部蛋白质的百分之几。从不到 5 个到几千个，差了约一千倍。

再来算一个让人心里一动的数字。阻遏蛋白在每个细菌里只有大约 10 个拷贝。细菌的体积约 1 立方微米，也就是 $10^{-15}$ 升。10 个分子的浓度是多少？
\[
\frac{10}{6.02\times 10^{23}\ \mathrm{mol^{-1}}\times 10^{-15}\ \mathrm{L}}\approx 1.7\times 10^{-8}\ \mathrm{mol/L}
\]
十亿分之十七摩尔每升。这么稀的浓度还能精确地把住开关，靠的是阻遏蛋白对操纵序列那段特定序列极强的亲和力（第 48 章的分子识别在这里发力）。但 10 个分子也意味着：它掉下来的那一两秒，全凭运气。同样的两个细菌，吃同样的乳糖，开灯的时刻可以差出好几分钟——这就是“随机波动”的根源，后面还会回来讲它。
\end{qiaoliang}

调节还不止发生在线路的输入端，输出也会绕回来捏住输入，这就是\textbf{反馈回路}。大肠杆菌合成色氨酸的线路是个经典例子：色氨酸本身是这条生产线末端的产品，可它一旦多起来，就会反过来当阻遏蛋白的“帮手”，把生产线自己的开关按住——产品多了就少生产，少了就多生产。你在第 31 章见过的化学平衡会自动朝抵消扰动的方向移动，这里的逻辑神似，只不过执行者是基因线路。正反馈也存在：有的基因产物会促进自己继续表达，像麦克风的啸叫一样把自己锁在“开”的状态——一个暂时的信号就能被记住很久。

一层层的开关、旋钮、反馈接在一起，就是\textbf{基因调控网络}。它不像梯子，更像一张互相指来指去的地铁图。这张网才是“基因如何工作”的完整答案，而操纵子只是其中一个画得最清楚的站。

\section{给基因线路画“电路图”}

讲了这么多口头描述，该让符号登场了。工程师画电路图，生物学家画调控网络图，约定惊人地相似：基因画成方框或横线，调节关系画成箭头——

\begin{itemize}[itemsep=2pt]
\item 箭头 $\rightarrow$：促进。$A \rightarrow B$ 表示 A 蛋白让 B 基因更亮。
\item 平头线 $\dashv$：抑制。$A \dashv B$ 表示 A 蛋白把 B 基因调暗。
\end{itemize}

用这个记号，乳糖操纵子的核心就一句话：
\[
\text{阻遏蛋白} \dashv \text{乳糖操纵子}, \qquad \text{乳糖} \dashv \text{阻遏蛋白}
\]
乳糖抑制了“抑制者”，负负得正，所以乳糖的最终效果是开灯。色氨酸线路则是：
\[
\text{色氨酸} \rightarrow \text{阻遏蛋白活性} \dashv \text{色氨酸合成基因}
\]
一个抑制环——产品掐住自己的生产线，负反馈。而胚胎发育里常见的“一旦开始就停不下来的命运决定”，画出来往往是一个基因指回自己的箭头：$\text{A} \rightarrow \text{A}$，正反馈。箭头本身不带数值，只说明“谁影响谁、方向是什么”；真正的强度要看结合的概率，也就是上一节的调光旋钮。别小看这两根线：现代系统生物学就是把成千上万个这样的箭头输入计算机，预测细胞会怎样行动。

\section{“两段生长”引出的诺贝尔奖}

\begin{zhengju}
这套理论不是哪个天才拍脑袋想出来的，它的起点是一条奇怪的曲线。

20 世纪 40 年代，法国巴斯德研究所的雅克·莫诺在研究细菌生长。他把细菌数量随时间画成曲线，大多数时候是一条平滑上升的线。可有一次，培养液里同时放了葡萄糖和乳糖，曲线变成了奇怪的形状：上升——走平——再上升，像两级台阶，中间隔着一段“停顿”。别人可能会把这段停顿当成实验误差抹掉，莫诺偏要问：细菌在那几十分钟里干了什么？

他先排除了一个解释：停顿不是因为细菌死了一批再长新的一批——显微镜下细胞数目没有掉。那么一定是细菌内部换了“装备”：吃葡萄糖用的酶原本就有，吃乳糖用的酶要现造。莫诺管这个现象叫“二次生长”，并开始追问：细胞怎么“知道”乳糖来了？

真正难的在后面。想知道开关的构造，需要改变开关本身看结果——这就是遗传学的老办法（第 63 章以来的思路）。莫诺和同事弗朗索瓦·雅各布筛选了一大批“开关失灵”的突变细菌：有的菌株不管有没有乳糖都拼命造乳糖酶（开关卡死在“开”），有的菌株永远造不出来（卡死在“关”）。1961 年，他们提出了操纵子模型：存在一个“阻遏物”，平时压住基因，乳糖来了就把它解除。“卡死在开”的突变，正是阻遏物坏掉了；把正常的阻遏物基因补进去，开关就恢复正常。

但故事到这里还只是\textbf{模型}——“阻遏物”是他们从开关的行为反推出来的东西，谁也没见过。怀疑者可以说：也许根本没有什么阻遏蛋白，也许有别的原因。科学上回应这种怀疑只有一个办法：把它抓出来。1966 年，吉尔伯特和穆勒-希尔真的把这种蛋白质从细菌里分离出来，证明它会专门抱住操纵序列那段 DNA，而且乳糖衍生物一来就松手——和模型预言的一模一样。模型里的“幽灵”变成了离心管里称得出质量的分子，操纵子学说这才算站稳。莫诺和雅各布因此获得了 1965 年诺贝尔奖（颁奖时蛋白分离工作刚刚完成，评委会其实冒了点险，赌对了）。

这个故事里最值得带走的方法论是：先看现象里“不该存在的台阶”，再用突变体把假想开关的每个零件逐一敲坏验证，最后把分子抓出来对质。三层证据凑齐之前，“阻遏物”始终只是个有待证实的假说。
\end{zhengju}

\begin{shiyan}
\textbf{观察酵母的“换档”}（在家可做）

材料：干酵母一小包、白砂糖、两个相同的透明矿泉水瓶、两个气球、温水（手感温热不烫，约 $35\,^{\circ}\mathrm{C}$，超过 $45\,^{\circ}\mathrm{C}$ 会烫死酵母）、勺子、记号笔。

步骤：
\begin{enumerate}[itemsep=1pt]
\item 两个瓶子各装半瓶温水。一瓶加两勺糖摇匀，另一瓶不加糖，作为对照。
\item 每瓶各倒入半勺干酵母，轻轻摇晃。
\item 把气球套在瓶口上，用记号笔在瓶身记下时间和初始状态。
\item 放在温暖处，每 10 分钟观察一次，持续一小时以上。记录：气球什么时候开始鼓？哪瓶先鼓？加糖那瓶是“立刻全速”还是“慢慢加速”？
\end{enumerate}

观察点：酵母遇到糖后，要把发酵生产线开足马力需要一点时间，气球往往是先慢后快地鼓起来——这和莫诺看到“台阶”是同一类现象：细胞换食谱时要现造装备。不加糖的对照瓶帮你确认：鼓气靠的是糖，而不是酵母遇水就产气。

安全提示：只用温水，不用热水瓶直接对嘴试温；实验后液体倒下水道、洗净瓶子即可；不要密封玻璃容器，气体积聚会撑破。
\end{shiyan}

\section{这套模型在哪里会失灵}

乳糖操纵子干净漂亮，但它是\textbf{细菌}的开关。把它原样搬到人身上，会连撞三堵墙。

第一堵墙：真核细胞（你、月季、果蝇）没有操纵子这种“一个总闸带一串基因”的排布，基因大多各有各的启动子，各自配一堆转录因子，调控比细菌繁复得多。增强子可以离它管的基因隔着几十万个字母，靠 DNA 折叠迂回接触——细菌的“路障挡门”图景在这里完全不适用。

第二堵墙：你的 DNA 不是裸露的，而是缠在组蛋白线轴上、再层层折叠压缩的染色质（第 66 章）。一段 DNA 如果被裹得严严实实，聚合酶连靠近都做不到——所以真核细胞的“开关”很大程度上是\textbf{这段 DNA 露出多少}的问题。细胞会把暂时不用的基因区域打包封存，把要用的区域松开来。更有意思的是，这种“打包方式”有时能在细胞分裂时被抄给下一代细胞，甚至（在某些动植物里，证据仍在争论中）传给后代——序列一个字母没改，使用方式却变了。这是第 74 章表观遗传的主题，这里只记住一句：调控可以分层，有的层比序列本身更“软”，却同样能被继承。

第三堵墙：就算是细菌，操纵子也远不是全貌。真实的调控网络里，乳糖线路的每个零件又连着别的线路，牵一发动全身。“一个开关控制一个功能”是教学模型，不是细胞的日常。

但模型失灵不等于没用。恰恰相反，正因为操纵子结构简单、开关分明，它被工程师改造成了生物技术里最好用的“零件”：把人的胰岛素基因接在乳糖操纵子的启动子后面装进大肠杆菌，平时关着灯让细菌安心繁殖，等细菌养够了量，往发酵罐里加诱导物，整罐细菌一起开灯造胰岛素——糖尿病人用的重组胰岛素，很大一部分就是这样生产的（第 55 章讲过胰岛素在代谢里的角色，第 83 章还会细讲它的工业化生产）。一个 1961 年的基础发现，如今每天都在药厂里开工。

\begin{wuqu}
“细胞只用它需要的基因，用不到的基因会慢慢退化消失。”——两处分不清。第一，在\textbf{个体的细胞里}，不用的基因只是被关掉了，序列完好无损：月季枝条上的细胞“本来”不长根，扦插后照样能打开长根的程序；植物的一小块组织甚至能在培养皿里重新长成完整植株。克隆羊多莉更直接：它来自一个成年羊乳腺细胞的细胞核——那个核里“关着”的制造整只羊的全部程序，被重新打开并完整执行了一遍。灯关了，不等于灯被拆了。第二，在\textbf{进化的漫长时间里}，某个性状彻底不用了（比如洞穴鱼的眼睛），相关基因确实可能因突变积累而损坏——但那是自然选择放松了维护的结果，以百万年计，跟单个细胞“关灯”完全不是一回事，更不是细胞“觉得没用就丢掉”。混淆这两层，就会把调节和突变、个体发育和进化搅成一锅粥（第 72 章讲突变，第 76 章起讲进化，会再遇到这对区别）。
\end{wuqu}

\section{回到那根划破的手指}

现在可以回答开场的问题了。伤口边缘的细胞里，制造肝、制造眼角膜的基因一个都不少——它们被层层关住：转录因子不在场，增强子沉默，DNA 区域打包封存。而“皮肤细胞身份”则被一组正反馈回路锁死：几个关键转录因子各自促进自己的表达，互相成就，让细胞“记得”自己是皮肤细胞。伤口愈合时打开的是修补程序——增殖、迁移、分泌胶原蛋白的相关基因被临时调高；长好之后，负反馈信号把增殖的灯关掉，刹车及时踩下（刹车失灵、灯关不掉，正是第 59 章谈过的癌症的一种来源）。

萤火虫屁股发光，是它的荧光素酶基因恰好在尾部那一节细胞里被打开；月季插条长根，是伤口激素信号拨动了一批转录因子，把“根”的亮度模式调了出来；你晒太阳变黑，是紫外线信号顺着一条蛋白质接力，摸到了黑色素细胞里那几个基因的旋钮。全是同一套逻辑：序列是书，调控是读法；书大家都一样，读法千变万化。

还有一个更深的味道值得咂摸：这套系统里没有一个“总司令”。开灯关灯，全是分子碰撞、结合、掉落这些局部事件的统计结果。可当成千上万个旋钮按网络图连在一起，整体行为却惊人地可靠——皮肤细胞几乎从不错长成肝细胞。可靠性不是来自某个指挥者，而是来自冗余：好几道闸、好几层反馈，坏了一道还有别的兜底。这和第一册反复出现的主题一脉相承：宏观的确定，从微观的概率里涌现。

\section{再往下挖：灯可以忽明忽暗}

\begin{tiaozhan}
还记得每个细菌里只有约 10 个阻遏蛋白分子吗？这意味着开关状态有真正的\textbf{随机性}：同一培养液里、基因完全相同的两个细菌，在某一瞬间，一个的操纵子可能恰好空着（灯亮），另一个的被占着（灯灭）。这叫基因表达的“噪声”。

噪声通常被细胞当作要压低的麻烦，但有时它会被\textbf{利用}。研究者在枯草杆菌里发现：环境恶劣时，少数个体会因为调控线路里的随机波动，碰巧跨过一个正反馈的门槛，进入“感受态”——一种能从外界吸收 DNA 的特殊状态。整个群体像买了一批彩票，随机中奖的少数个体获得了别的生存选项。更惊人的是，发育中的胚胎也有类似情节：某些组织里，细胞本来完全相同，随机波动先把某个细胞推过门槛，正反馈立刻把它锁进新状态，它再向邻居发出“你别变成我这样”的抑制信号——一盏随机亮起的灯，最终决定了谁当神经元、谁当表皮。命运的图纸不全是写死的，有一部分是掷骰子加自锁定。这部分内容涉及反馈回路的数学（双稳态、分岔），跳过不影响主线，但如果你好奇“随机怎么变成可靠的模式”，这是当代发育生物学最活跃的前沿之一。
\end{tiaozhan}

说到这里，一个更根本的问题已经绕不开了：我们一直在说“基因的表达”“基因被关掉”，可“基因”这个词本身指的是什么？是“制造蛋白质的段落”？那制造调控 RNA、不制造蛋白质的段落算不算？调控序列、增强子算不算基因的一部分？一段 DNA 被剪接成好几种产物，算一个基因还是几个？你可能会惊讶：这个中学生物课本里张口就来的词，到今天都没有一个让所有人满意的定义。下一章，我们就来给“基因”这个词本身动一次手术。

\section*{练一练}

\begin{sikaoti}
\item 用文字或符号（$\rightarrow$ 与 $\dashv$）画出乳糖操纵子的信息流图：从“培养液里出现乳糖”到“乳糖酶浓度上升”，标出每一环是促进还是抑制。再用同一套记号画出色氨酸合成的负反馈环。
\item 阻遏蛋白从操纵序列上掉下来是概率事件。假设没有乳糖时它 99\% 的时间坐在 DNA 上，有足量乳糖时只有 1\% 的时间坐在上面。定性说明：为什么中间浓度（比如 50\%）对应的是“半亮的灯”，而不是“一半细胞全开、一半全关”？（提示：想想一个开关在几分钟内会切换多少次。）
\item 莫诺的二次生长曲线中间有一段“停顿”。请提出两个不同的假说来解释停顿（一个已被他的观察排除），并说说你会设计什么观察来区分剩下的假说。
\item 同一株月季的两个细胞，一个在花瓣里、一个在根里，DNA 序列完全相同，制造的蛋白质却大半不同。请画一幅“物质流/信息流”示意：从“相同的 DNA 书”出发，经过哪些环节，分叉出两种不同的蛋白质清单？在分叉点上标注“是谁在读哪一页”。
\item 生物技术公司用改造过的乳糖操纵子控制细菌生产胰岛素。为什么工程师要故意把灯“平时关着”，而不是让细菌一直全速生产？从细菌的能量账本（第 50 章）想想常开的代价。
\item 洞穴鱼的眼睛退化常被误引为“细胞不用的基因会消失”的证据。根据本章内容，指出这个误引混淆了哪两个时间尺度、哪两个层面的过程？
\end{sikaoti}
